\section{总结与延伸}

\subsection{核心要点回顾}

DDIM 是扩散模型加速采样的奠基性工作，其核心贡献可以归纳为以下几点：

\begin{enumerate}
    \item \textbf{非马尔可夫前向过程}：通过直接定义跳跃分布和后验分布（而非前向转移分布），构建了比 DDPM 更通用的扩散框架。后验分布中引入了自由参数 $\sigma_t$，控制采样的随机性程度。

    \item \textbf{ELBO 不变性}：由于跳跃分布与 DDPM 保持一致，DDIM 的训练目标与 DDPM 完全等价。这意味着已有的 DDPM 模型可以直接用于 DDIM 采样，无需任何重新训练。

    \item \textbf{确定性采样}：当 $\sigma_t = 0$ 时，采样过程变为完全确定性的，建立了噪声空间到数据空间的双射映射。这带来了可复现性、语义隐空间、以及与 ODE 求解器的联系。

    \item \textbf{加速生成}：通过子序列采样策略，可以将步数从 1000 降低到约 50，实现 20 倍加速。确定性采样模式在步数减少时表现出更强的鲁棒性。

    \item \textbf{隐空间操作}：确定性 DDIM 的编码-解码对称性使得隐空间插值等语义操作成为可能。
\end{enumerate}

\begin{importantbox}{DDIM 的深远影响}
DDIM 不仅本身是一个实用的加速采样方法，更重要的是它揭示了扩散模型一个关键的设计原则：\textbf{训练与推理的解耦}。模型训练一次后，可以在推理时自由选择不同的采样策略（步数、随机性程度、ODE 求解器等）。这一原则被后续的几乎所有扩散模型加速工作所继承，包括 DPM-Solver、Consistency Models、Progressive Distillation 等。
\end{importantbox}

\subsection{后续发展方向}

基于 DDIM 的框架，后续研究从多个方向进一步推进了扩散模型的采样效率：

\begin{enumerate}
    \item \textbf{高阶 ODE 求解器}（DPM-Solver, DPM-Solver++）：将确定性 DDIM 视为概率流 ODE 的 Euler 方法，使用更高阶（2阶、3阶）的求解器可以在相同步数下获得更高精度，或在更少步数下达到相同质量
    \item \textbf{连续时间 SDE/ODE 框架}（Score SDE）：将离散时间的 DDPM/DDIM 推广到连续时间，统一了 score-based models 和扩散模型
    \item \textbf{知识蒸馏}（Progressive Distillation, Consistency Distillation）：通过蒸馏进一步将步数降低到 1-4 步
    \item \textbf{Flow Matching}：从最优传输的角度重新设计扩散路径，实现更直的采样轨迹和更少的步数
\end{enumerate}

\subsection{拓展阅读}

\begin{itemize}
    \item Song, J., Meng, C., \& Ermon, S. (2021). Denoising Diffusion Implicit Models. \textit{ICLR 2021}. \href{https://arxiv.org/abs/2010.02502}{arXiv:2010.02502}
    \item Song, Y., Sohl-Dickstein, J., Kingma, D. P., Kumar, A., Ermon, S., \& Poole, B. (2021). Score-Based Generative Modeling through Stochastic Differential Equations. \textit{ICLR 2021}. \href{https://arxiv.org/abs/2011.13456}{arXiv:2011.13456}
    \item Ho, J., Jain, A., \& Abbeel, P. (2020). Denoising Diffusion Probabilistic Models. \textit{NeurIPS 2020}. \href{https://arxiv.org/abs/2006.11239}{arXiv:2006.11239}
    \item Lu, C., Zhou, Y., Bao, F., Chen, J., Li, C., \& Zhu, J. (2022). DPM-Solver: A Fast ODE Solver for Diffusion Probabilistic Model Sampling. \textit{NeurIPS 2022}. \href{https://arxiv.org/abs/2206.00927}{arXiv:2206.00927}
    \item Song, Y., \& Dhariwal, P. (2023). Consistency Models. \textit{ICML 2023}. \href{https://arxiv.org/abs/2303.01469}{arXiv:2303.01469}
\end{itemize}

%% --- Fin del contenido --- %%

\end{document}
